% ---------------------------------------------------------------------------
\appendix
\renewcommand{\thesection}{\Alph{section}}

\section{Experiments behind the main findings}

Experiment codes refer to the repository (files \texttt{preregistrazioni/}, \texttt{esperimenti/}, \texttt{risultati/}).

\begin{center}\footnotesize
\begin{tabularx}{\linewidth}{>{\raggedright\arraybackslash}p{5.4cm}Yp{3.1cm}}
\toprule
finding & main experiments & replications \\
\midrule
Statistical fingerprint, isolation from languages & e01--e06, e13 & IT, GC \\
Weak and predictable space, cuts yielding words, vocabulary filling the forms & e3a49, e3a55, e3a58, e3a61, e3a67,
e3a99--e3b05 & e3a50, e3a56, e3a60, e3a68 \\
Junction and junction rules & e377, e380, e390, e3a01--e3a03, e3a41 & e388, e395, e3a04, e3a05, e3a34 \\
Closed line, drawing gap, line-edge forms & e384, e386, e387, e389, e3a09 & e395, e3a43 \\
Copying from the line immediately above & e340, e343, e347, e372, e385, e392, e394, e396, e3b28 & e3a80, e3a84 \\
Left-margin avoidance & e3a25--e3a27, e3a33, e3a35, e3a44 & e3a33 (IT, languages A and B) \\
Paragraph, page, bifolium, quire, languages A and B & e302, e307--e316, e321, e350, e356, e378 & e3a15, e3a21 \\
Immediate repetition and its shape & e02, e3a86--e3a90, e3b25 & e3a88, e3a91, e3a92 \\
Short state of the choices & e3c33, e3c48, e3c52, e3c55, e3c57, e3c66, e3c72, e3c81 & IT, GC (e3c51), hands, sections
(e3c54) \\
Slow state & e3c73--e3c75 & e3c76 (scribes) \\
Drift along the line & e3c13, e3c17, e3c20--e3c24 & IT \\
Scribe battery & e3c50, e3c58--e3c65, e3c76--e3c80 & --- \\
Copiale & e3c61 & --- \\
Decipherments and readings excluded & e14, e16--e20, e26, e27, e133, e163, e163b, e212b, e216, e219, e222, e223, e249,
e269 & positive controls in every test \\
Historical cipher mechanisms & e3c68--e3c71 & synthetic controls in every test \\
Constructed languages, gibberish, generators & e346, e374, e377, e382, e399, e3a95, e3c37, e3c49, e3c53 & --- \\
\bottomrule
\end{tabularx}
\end{center}

% ---------------------------------------------------------------------------
\section{Internal checks: provisional indications corrected by the method}

None of the results of this work had been published before. During the work, however, some intermediate measurements gave
indications that later checks corrected. We report them because they show which measurements are fragile and how the
controls corrected them: anyone redoing the work can avoid the same missteps, and readers can judge how far the final
conclusions have been tested.

\paragraph{Corrections due to the bias of the agreement measure (Section~3.3)}
\begin{enumerate}
\item A first measurement attributed to the Anglo-Saxon scribe of the \emph{Hatton Gospels} the same ``window'' as the
  Voynich. For that choice the measure was inflated ($+0.25$); corrected, the scribe has agreement between adjacent words
  but not at two-three words.
\item Volapük seemed to have an agreement similar to the Voynich; with the corrected measure it is zero.
\item The shape of the Voynich's agreement (how long it lasts at two-three words) seemed outside the range of languages;
  corrected, it lies at the top of the range, not outside.
\item An observation about ``neutral'' words between two words making the same choice suggested that the state travels only
  on occurrences of the choice; it was an effect of the same bias.
\end{enumerate}

\paragraph{Descriptions made more precise by later measurements}
\begin{enumerate}[resume]
\item The state of the choices looked like a ``window'' of three words then dropping to zero; instead it fades gradually
  (half-life about three words) and there is in addition a slow component carrying from one line to the next.
\item The hypothesis that each line has its own ``setting'' of the choices (like a key per line) did not hold: agreement
  carries from one line to the next.
\item After the test on Glen Claston's form-only variants we had read the state as tied to ``which word is written, not
  how the letter is drawn''; a real scribe with a window in a form-only choice showed that the generalisation was too broad
  (and his window, in turn, comes from similar words).
\item An uncorrected measure indicated that hand 1 lacked the state of the choices; with the corrected measure all hands
  have it equally.
\item A measure indicated that the memory of the choices resets at the line break; it was an effect of comparing with the
  page average (line ends and starts have their own forms).
\item The first left-margin test used a defective comparison; with the correct comparison the result holds.
\item A measure attributed two thirds of the small link between whole words to the paragraph; it was a mechanical effect
  of smaller groups.
\item Rare forms seemed to be copied from the line above more than others; with the second transliteration only a trend
  remains.
\end{enumerate}

\paragraph{Formally positive outcomes that were artefacts}
\begin{enumerate}[resume]
\item A ``one word = one syllable'' reading candidate and a German candidate in a content-guided search: the first came out
  identical on the shuffled Voynich, the second changed sign when only the starting points of the search were changed.
\end{enumerate}

\paragraph{Criteria phrased too sharply}
\begin{enumerate}[resume]
\item In three cases the criterion decided in advance called ``absent'' or ``reset'' an effect whose uncertainty interval
  touched zero (the state in the herbal, at the drawing gap, in some sections). In all three the estimate was positive and
  similar to the rest: these were insufficient data, not absence. Hence the rule, which we suggest to others too, of
  saying ``absent'' only when the interval lies entirely below the reference value.
\end{enumerate}

% ---------------------------------------------------------------------------
\section{Glossary}

\begin{description}[leftmargin=0pt,itemsep=2pt]
\item[EVA] conventional alphabet to transliterate the Voynich into Latin letters; it does not indicate sounds.
\item[glyph] a sign of the manuscript; \eva{ch}, \eva{sh}, the gallows and their compounds count as one.
\item[gallows] the tall signs \eva{k}, \eva{t}, \eva{p}, \eva{f} of EVA.
\item[Currier languages A and B] the two varieties of writing into which the pages divide.
\item[junction] how much information the last glyph of a word gives about the first glyph of the next, beyond chance.
\item[junction rule] a rule linking the form of a word to that of its neighbour (sandhi-like).
\item[closed line] the junction does not carry over the line break.
\item[state of the choices] the tendency to repeat the same variant for a few words.
\item[in-line drift] the decline of marked variants from left to right.
\item[preregistration] a written description, before running, of question, measure and criterion.
\item[positive / negative control] a text in which the answer is / is not present.
\item[homophonic cipher] the same letter can be written with several symbols.
\item[verbose cipher] each letter becomes a group of several glyphs.
\item[self-citation] a procedure in which the writer copies and modifies words already written.
\item[facsimile-level transcription] a transcription preserving the letter forms of the manuscript.
\item[95\% interval] the range in which the true value lies, with good confidence.
\end{description}

% ---------------------------------------------------------------------------
\section{Sources and licences}

\subsection{Data}

\paragraph{The Voynich}
\begin{itemize}
\item R. Zandbergen, transliteration files of the Voynich manuscript in IVTFF format,
  \url{https://www.voynich.nu/transcr.html}: Zandbergen--Landini (ZL3b, EVA), Takahashi (IT2a), Glen Claston (GC2a, v101
  alphabet); format described in R. Zandbergen, \emph{IVTFF -- Intermediate Voynich MS Transliteration File Format},
  v2.0.1 (2025).
\item Radiocarbon dating of the vellum (1404--1438, 95\%): NSF-Arizona AMS Laboratory, University of Arizona, team led by
  G. Hodgins, four samples measured in 2009; University of Arizona press release, 10 February 2011; account by R.
  Zandbergen, \url{https://www.voynich.nu/extra/carbon.html}.
\end{itemize}

\paragraph{Comparison texts}
\begin{itemize}
\item C. Christodouloupoulos and M. Steedman (2015), ``A massively parallel corpus: the Bible in 100 languages'',
  \emph{Language Resources and Evaluation} 49(2): 375--395, doi:\url{10.1007/s10579-014-9287-y} (data: \url{christos-c/bible-corpus},
  CC0).
\item The Latin Library (collection \url{cltk/lat_text_latin_library}); Roman breviary (Divinum Officium, MIT).
\item D. E. Gaskell and C. L. Bowern (2022), ``Gibberish after all? Voynichese is statistically similar to human-produced
  samples of meaningless text'', in C. Layfield and J. Abela (eds.), \emph{Proceedings of the 1st International Conference
  on the Voynich Manuscript 2022 (VOY2022)}, CEUR Workshop Proceedings 3313, paper 4 (2023),
  \url{https://ceur-ws.org/Vol-3313/paper4.pdf}. The meaningless texts were handwritten by volunteers (42 in the study; 38
  texts in the published data); data licence: modified MIT.
\item M. A. Greshko (2025), ``The Naibbe cipher: a substitution cipher that encrypts Latin and Italian as Voynich
  Manuscript-like ciphertext'', \emph{Cryptologia}, published online 26 November 2025, doi:\url{10.1080/01611194.2025.2566408}
  (code: \url{greshko/naibbe-cipher}).
\item T. Timm and A. Schinner (2020), ``A possible generating algorithm of the Voynich manuscript'', \emph{Cryptologia}
  44(1): 1--19, doi:\url{10.1080/01611194.2019.1596999} (code: \url{TorstenTimm/SelfCitationTextgenerator}, MIT).
\item \textbf{Menota}, Medieval Nordic Text Archive, \url{https://www.menota.org} (catalogue on Clarino, University of
  Bergen). The licence is stated text by text; for the six texts used (AM 519 a 4to, AM 677 4to, AM 60 4to, AM 242 fol,
  Holm A 10, AM 302 fol) it is CC BY-SA 4.0. Each text should be cited with the data of its header (editor, version).
\item \textbf{ReF}: K.-P. Wegera, H.-J. Solms, U. Demske, S. Dipper (2021), \emph{Reference Corpus of Early New High
  German (1350--1650)}, version 1.0.2, Zenodo, doi:\url{10.5281/zenodo.5793616}. Zenodo states CC BY 4.0, the LICENSE file inside
  the archive CC BY-SA 4.0: we follow the more restrictive.
\item \textbf{Copiale}: K. Knight, B. Megyesi, C. Schaefer (2011), ``The Copiale Cipher'', in \emph{Proceedings of the 4th
  Workshop on Building and Using Comparable Corpora (BUCC)}, ACL, pp. 2--9, \url{https://aclanthology.org/W11-1202/}; and K.
  Knight, B. Megyesi, C. Schaefer, ``The Secrets of the Copiale Cipher'', \emph{Journal for Research into Freemasonry and
  Fraternalism} 2(2): 314--324 (2011 volume, published 2012), doi:\url{10.1558/jrff.v2i2.314}. Transcription and decipherment
  from the project page of Stockholm University. Dating: mid-eighteenth century (the inscription ``1866'' in the manuscript
  is an owner's mark).
\end{itemize}

\subsection{Readings and decipherments tested}
\begin{itemize}
\item S. Bax (2014), \emph{A proposed partial decoding of the Voynich script}, self-published PDF, January 2014 (read from
  the web archive).
\item S. B. Vatne (2021), \emph{Cracking the Voynich Cipher}, self-published PDF, October 2021.
\item G. Cheshire (2019), ``The Language and Writing System of MS408 (Voynich) Explained'', \emph{Romance Studies} 37(1):
  30--67, doi:\url{10.1080/02639904.2019.1599566}.
\item S. Schechter (2026), \emph{voynich-decoded}, GitHub repository,
  \url{https://github.com/scott-schechter/voynich-decoded} (no licence stated; used as a local copy, not redistributed).
\item A. Gatta (2026), \emph{voynich-toolkit} (software), Zenodo, doi:\url{10.5281/zenodo.19226178} (concept DOI; the version
  tested is v0.26.0), MIT licence.
\end{itemize}

\subsection{Literature}
\begin{itemize}
\item P. H. Currier (1976), ``Papers on the Voynich Manuscript'', in M. E. D'Imperio (ed.), \emph{New Research on the
  Voynich Manuscript: Proceedings of a Seminar}, Washington (typescript); transcription by J. Guy and J. Reeds (1992),
  \url{https://www.voynich.nu/extra/curr_main.html}.
\item L. Fagin Davis (2020), ``How Many Glyphs and How Many Scribes? Digital Paleography and the Voynich Manuscript'',
  \emph{Manuscript Studies} 5(1): 164--180, doi:\url{10.1353/mns.2020.0011} (the five hands used here).
\item M. A. Montemurro and D. H. Zanette (2013), ``Keywords and co-occurrence patterns in the Voynich manuscript: an
  information-theoretic analysis'', \emph{PLoS ONE} 8(6): e66344, doi:\url{10.1371/journal.pone.0066344}.
\item B. Hauer and G. Kondrak (2016), ``Decoding Anagrammed Texts Written in an Unknown Language and Script'',
  \emph{Transactions of the Association for Computational Linguistics} 4: 75--86, doi:\url{10.1162/tacl_a_00084}.
\item L. Rozanova and A. Temerev (2026), ``A Glyph Is Not a Letter, a Token Is Not a Word, a Space Is Not a Space: What
  the Units of Voynichese Are Not'', arXiv:2608.17096.
\item L. B. Alberti, \emph{De componendis cifris} (c.~1466); J. Trithemius, \emph{Polygraphiae libri sex} (1518); F.
  Bacon, \emph{De dignitate et augmentis scientiarum} (1623), with the description of the biliteral cipher.
\end{itemize}

\noindent\textit{All references were checked on publisher, repository or archive pages (October 2026). For Bax the title
page of the PDF could not be re-read; for the radiocarbon dating the source is a university press release, not a
scientific article.}
